\section{Explicit ODEs derivation}

\label{2305:app:square_equations}
\subsection{Derivation of the process}
Let's start by reminding the definition of \emph{displacement} at step \(\i\)
\begin{equation}
\mathcal{E}^{\i} \coloneqq (\lambda_{\star}^\i)^2 - \frac{1}{p}\sum_{j=1}^{\hids}a_{j}(\lambda_{j}^{\nu})^2 + \sqrt{\noise} z^{\i},
\end{equation}
from which it's easy to write the loss function
\begin{equation}
    \ell{(y^{\nu},f_{\Theta^{\nu}}(x^{\nu}))} = \frac12 \left(\mathcal{E}^{\i}\right)^2.
\end{equation}
The gradient respect to the parameters is given
\begin{equation}\begin{split}
    \partial_{a_j} \ell{(y^{\nu},f_{\Theta^{\nu}}(x^{\nu}))} &= -\frac1p \mathcal{E}^{\i} (\lambda_{j}^{\nu})^2 \\
    \nabla_{\w_j} \ell{(y^{\nu},f_{\Theta^{\nu}}(x^{\nu}))}  &= -\frac1p \mathcal{E}^{\i} 2 a^\i_j \lambda_{j}^{\nu} x^\i
\end{split}\end{equation}
Using \(\gamma\) as learning rate for the weights \(\w_j\) and \(\sfrac{\gamma}{d}\) for the second layer, we have the following update equations
\begin{equation}\begin{split}
    a^{\i+1}_j &= a^{\i}_j + \frac{\gamma}{pd} \mathcal{E}^{\i} (\lambda_{j}^{\nu})^2 \\
    \w_j^{\i+1} &= \w_j^{\i} + \frac{\gamma}{p} \mathcal{E}^{\i} 2 a^\i_j \lambda_{j}^{\nu} x^\i
\end{split}\end{equation}
Applying the definition of the sufficient statistics, \(m_j = \sfrac{\w_j\w_\star}{d}\) and \(\Q_{jl} = \sfrac{\w_j\w_l}{d}\), we can recover Equations~\eqref{2305:eq:def:overlap_process}.

\subsection{Explicit ODE}
To get the explicit form of our ODEs we need to compute some expected value over the preactivations \((\lambda_\star,\lambda)\).
These are gaussian variables, whose correlation matrix is given by \(\Omega\).
Similarly, the population risk is defined as
\[
    \risk = \mathbb E_{(\lf,\lf^{\star})\sim\mathcal{N}(0_{\hids+1}, \Omega)}\left[\frac12 \mathcal{E}^2\right].
\]

Let's look close to the random variable we need for this expected values, expressing them just as function of local fields.
To be more concise, from now on with \(\mathbb E\) we always mean the expected value over \((\lf,\lf^{\star})\sim\mathcal{N}(0_{\hids+1}, \Omega)\).
For the risk we just need
\[\begin{split}
  \EV{\mathcal{E}^2} =& \lambda_\star^4 + \frac{1}{p^2} \sum_{j, l =1}^p a_j a_l\EV{\lf_j^2 \lf_l^2} - \frac{2}{p} \sum_{j=1}^p a_j\EV{\lf_j^2 \lf_\star^2},
\end{split}\]
while for the ODE of \(\bar{a}\) 
\[\begin{split}
    \EV{\mathcal{E} \lf_j^2} = \EV{\lambda_{\star}^2\lf_j^2} - \frac{1}{p}\sum_{l=1}^{\hids}a_{l}\EV{\lambda_{l}^2\lf_j^2},
\end{split}\]
where we omitted the noise part since it averages out with the expectation.
The equation for \(\bar{m}\) requires 
\[\begin{split}
    2a_j\EV{\mathcal{E} \lf_j\lf_\star} = 2\EV{\lambda_{\star}^2\lf_j\lf_\star} - \frac{2}{p}\sum_{l=1}^{\hids}a_j a_{l}\EV{\lambda_{l}^2\lf_j\lf_\star},
\end{split}\]
while we need two different expectations for \(\bar\Q\)
\[\begin{split}
    2(a_j+a_l)\EV{\mathcal{E} \lf_j\lf_\star} &= 2(a_j+a_l)\left[\EV{\lambda_{\star}^2\lf_j\lf_l} - \frac{1}{p}\sum_{s=1}^{\hids}a_s\EV{\lambda_{s}^2\lf_j\lf_l}\right] \quad\text{and} \\
    %
    4\frac{\gamma}{p}{\mathcal{E}}^{2} a_{j}a_{l}\lambda_{j}\lambda_{l}  = 16\frac{\gamma}{p} a_{j}a_{l} \Bigg[ 
        &  \EV{\lf_j \lf_l\lf_\star^4}  - \frac{2}{p} \sum_{s =1}^{p} a_s \EV{\lf_j \lf_l \lf_\star^2 \lf_s^2} +\\
           &\qquad+ \frac{1}{p^2} \sum_{ s,  t =1}^{p}\left(a_sa_t\EV{\lf_j \lf_l \lf_s^2 \lf_t^2}+\Delta\EV{\lf_j \lf_l}\right) \Bigg].
\end{split}\]
We are left to compute some distribution moments of a multivariate Gaussian of second, fourth and sixh order.
In the \href{https://github.com/IdePHICS/EscapingMediocrity}{GitHub repository} can be found a Mathematica script to address this task;
alternatively Isserlis' Theorem can be applied.
We introduce a shorthand in the notation \[\omega_{\alpha\beta} \coloneqq \left[{\Omega}\right]_{\alpha\beta},\]
where the indices \(\alpha\) and \(\beta\) can discriminate between local fields \(\lambda\) (if \(\alpha,\beta \in [1,\dots,p]\), and \(\lambda_\star\) (if  \(\alpha,\beta = p+1\)). The final result are given by
\begin{align*}
  \EV{\lf_\alpha \lf_\beta} &= \omega_{\alpha\beta}\\
  %
  \EV{\lf_\alpha^2 \lf_\beta^2} 
  &= \omega_{\alpha\alpha}  \omega_{\beta\beta} + 2 \omega_{\alpha\beta}^2 \\
  %
  \EV{\lf_\alpha \lf_\beta \lf_\gamma^2}
  &= \omega_{\alpha\beta}\omega_{\gamma\gamma} + 2 \omega_{\alpha\gamma}\omega_{\beta\gamma}\\
  % 
  \EV{\lf_\alpha \lf_\beta \lf_\gamma^2 \lf_\delta^2}
  &= \omega_{\alpha\beta}\omega_{\gamma\gamma}\omega_{\delta\delta} + 2\omega_{\alpha\beta}\omega_{\gamma\delta}^2 + 2\omega_{\alpha\gamma}\omega_{\beta\gamma}\omega_{\delta\delta} +\\
  &\quad4\omega_{\alpha\gamma}\omega_{\beta\delta}\omega_{\gamma\delta}+4\omega_{\alpha\delta}\omega_{\beta\gamma}\omega_{\gamma\delta} + 2\omega_{\alpha\delta}\omega_{\beta\delta}\omega_{\gamma\gamma}
\end{align*}
By retracing all steps backward and making the necessary substitutions,
we can arrive at an explicit form of the ODEs and population risk. 
While the full risk expression can be found in Equation~\eqref{2305:eq:def:risk}, we report here just the case \(a_j=1\) for the ODEs since they have a compact matrix form
\begin{subequations}\label{2305:eq:full_squared_ode}\begin{align}
    \label{2305:eq:quadraticODE_M}
    \dod{\M}{t} &=
        2\left(\rho - \frac{\Tr{\Q}}p\right)\M +
        4\left(\rho\M - \frac{\Q\M}{p}\right) \\
    %
    \label{2305:eq:quadraticODE_Q}
    \begin{split}
        \dod{\Q}{t} &= 
            4\left(\rho - \frac{\Tr{\Q}}p\right)\Q +
            8\left(\frac{\M\M^\top}{p} - \frac{\Q^2}{p}\right) \\
        &\quad+\frac{4 \gamma}{p} \Biggr\{
            \left[3\rho^2\Q + 12\rho\M\M^\top\right] +\frac1{p^2}\left[\left(\Tr{\Q}^2+2\Tr{\Q^2}\right)\Q +
                            4\Tr{\Q}\Q^2 + 8\Q^3
                        \right]\\
            &\qquad\quad\quad
            -\frac2{p}\Biggl[\left(\rho\Tr{\Q}+2\Tr{\M\M^\top}\right)\Q +
                            2\Tr{\Q}\M\M^\top\\
                            &\qquad\qquad\qquad\quad+ 2\rho\Q^2+ 
                            4\left(\M\M^\top\Q + \Q\M\M^\top\right)
                        \Biggl] + \Delta \Q\Biggr\},
    \end{split}
\end{align}\end{subequations}
where \(\rho \coloneqq \sfrac{\w_\star^2}{d}\). 
For completeness, this is Equation~\eqref{2305:eq:def:risk} for the case \(a_j=1\)
\begin{equation} \label{2305:eq:risk_square}\begin{split}
  \mathcal{R}(\Omega)
       =  \frac{3+\Delta}{2}
           - \frac{\rho\Tr{\Q}+2\Tr{\M\M^\top}}{p}
           + \frac{\Tr{\Q}^2+2\Tr{\Q^2}}{2p^2}.
\end{split}\end{equation}
% %%%%%%%%%%%%%%%%%%%%%%%%%%%%%%%%%%%%%%%%%%%%%%%%%%%%%%%%%%%%%%%%%%%%%%%%%%%%%%%
